\documentclass[12pt,a4paper]{article}
\usepackage[utf8]{vietnam}
\usepackage{amsmath,amssymb}
\usepackage{tabularx}
\usepackage{tikz}
\usepackage{xcolor}
\usepackage[top=0.6cm, bottom=0.6cm, left=0.6cm, right=0.6cm]{geometry}
\usepackage[dethi]{ex_test}

\setlength{\parskip}{0pt}

% Dòng kẻ chấm mờ đều đặn chuẩn viết tay (khoảng cách 7.2mm rộng rãi)
\newcommand{\dongke}[1]{%
  \begingroup%
  \setlength{\parskip}{0pt}%
  \baselineskip=20.5pt%
  \vspace{1.5mm}%
  \foreach \i in {1,...,#1}{%
    \noindent{\color{gray!60}\dotfill}\par%
  }%
  \vspace{1mm}%
  \endgroup%
}

% Định nghĩa ô bài tập tự luận trong bảng
\newcommand{\oBaiTap}[3][6]{%
  \vspace{1.5mm}%
  \noindent\textbf{#2} #3\par%
  \dongke{#1}%
}

% Câu trắc nghiệm 1 dòng (4 đáp án, mỗi đáp án 0.25\linewidth)
\newcommand{\cauTN}[6]{% \cauTN{số câu}{đề}{A}{B}{C}{D}
  \noindent\textbf{Câu #1.} #2\par\vspace{0.4mm}%
  \noindent%
  \makebox[0.25\linewidth][l]{\textbf{A.} #3}%
  \makebox[0.25\linewidth][l]{\textbf{B.} #4}%
  \makebox[0.25\linewidth][l]{\textbf{C.} #5}%
  \makebox[0.25\linewidth][l]{\textbf{D.} #6}\par\vspace{0.8mm}%
}

% Câu trắc nghiệm 2 dòng (mỗi dòng 2 đáp án, mỗi đáp án 0.5\linewidth) cho đáp án dài
\newcommand{\cauTNB}[6]{% \cauTNB{số câu}{đề}{A}{B}{C}{D}
  \noindent\textbf{Câu #1.} #2\par\vspace{0.4mm}%
  \noindent%
  \makebox[0.5\linewidth][l]{\textbf{A.} #3}%
  \makebox[0.5\linewidth][l]{\textbf{B.} #4}\par\vspace{0.4mm}%
  \noindent%
  \makebox[0.5\linewidth][l]{\textbf{C.} #5}%
  \makebox[0.5\linewidth][l]{\textbf{D.} #6}\par\vspace{0.8mm}%
}

\begin{document}
\pagestyle{empty}

% =============================================================================
% ========================= TOÁN 6 - TUẦN 05 - BUỔI 02: MÃ ĐỀ 101 =============
% ========================= MẶT TRƯỚC (TRANG 1/2) ==============================
% =============================================================================

\noindent
\begin{tabularx}{\linewidth}{|X|c|}
\hline
\textbf{PHIẾU KIỂM TRA ĐỊNH KỲ TOÁN 6 - TUẦN 05 - BUỔI 02} & \textbf{Mã đề: 101 - Trang 1/2} \\
\textit{Chủ đề: Dấu hiệu chia hết (2, 3, 5, 9) --- Tìm ƯCLN --- Tứ giác thực tế} & \textit{Thời gian: 45 phút} \\
\hline
\end{tabularx}

\vspace{1mm}

\noindent
\begin{tabularx}{\linewidth}{|X|}
\hline
\rule[-3.5mm]{0pt}{10mm}\textbf{Họ và tên học sinh:} \hfill \textbf{Mã HS:} \makebox[3.5cm]{} \\
\hline
\end{tabularx}

\vspace{1mm}

\noindent
\begin{tabularx}{\linewidth}{|p{3.2cm}|X|}
\hline
\centering\textbf{ĐIỂM SỐ} & \centering\textbf{LỜI NHẬN XÉT CỦA THẦY CÔ} \tabularnewline
\hline
\rule{0pt}{1.3cm} & \tabularnewline
\hline
\end{tabularx}

\vspace{1.2mm}

\noindent\textbf{BẢNG ĐIỀN ĐÁP ÁN TRẮC NGHIỆM (Mỗi câu đúng 0,5 điểm):}
\vspace{0.8mm}

\noindent
\begingroup
\setlength{\tabcolsep}{3pt}
\renewcommand{\arraystretch}{1.35}
\begin{tabularx}{\linewidth}{|c|*{5}{>{\centering\arraybackslash}X|}}
\hline
\textbf{Câu} & \textbf{1} & \textbf{2} & \textbf{3} & \textbf{4} & \textbf{5} \\
\hline
\textbf{Đáp án} & & & & & \\
\hline
\end{tabularx}
\endgroup

\vspace{1.5mm}

\noindent\textbf{\large PHẦN I. TRẮC NGHIỆM KHÁCH QUAN (2,5 điểm)}
\vspace{1mm}

\cauTN{1}{Trong các số sau: $125, 230, 342, 515$, số chia hết cho cả $2$ và $5$ là:}{$125$}{$230$}{$342$}{$515$}

\cauTN{2}{Số nào sau đây chia hết cho $3$?}{$145$}{$253$}{$357$}{$418$}

\cauTN{3}{Cho số $A = \overline{35*}$. Để $A$ chia hết cho $9$ thì chữ số $*$ thích hợp là:}{$1$}{$4$}{$7$}{$9$}

\cauTNB{4}{Khẳng định nào sau đây là \textbf{sai}?}{Nếu $a \ \vdots \ 4$ và $b \ \vdots \ 4$ thì $(a+b) \ \vdots \ 4$}{Nếu $a \ \vdots \ 5$ và $b \ \not\vdots \ 5$ thì $(a+b) \ \not\vdots \ 5$}{Nếu $(a+b) \ \vdots \ 6$ thì mỗi số $a, b$ đều chia hết cho 6}{Nếu $a, b \ \vdots \ 3$ thì $(a-b) \ \vdots \ 3$ (với $a \ge b$)}

\cauTN{5}{Trong các số: $120, 235, 348, 540$, số chia hết cho cả $2, 3, 5$ và $9$ là:}{$120$}{$235$}{$348$}{$540$}

\vspace{1.2mm}

\noindent\textbf{\large PHẦN II. TỰ LUẬN (7,5 điểm)}

\noindent\textbf{Câu 1 (2,5 điểm). Thực hiện phép tính (hợp lý nếu có thể):}
\vspace{0.8mm}

\noindent
\begin{tabularx}{\linewidth}{|X|X|}
\hline
\oBaiTap[6]{a)}{$38 + 24 \cdot 2 - 2^4$} & 
\oBaiTap[6]{b)}{$43 \cdot 78 + 43 \cdot 22 - 800$} \\
\hline
\oBaiTap[6]{c)}{$120 : [54 - (50 - 2^3 \cdot 2)]$} & 
\oBaiTap[6]{d)}{$240 : \{ 180 - [75 + (8 - 3)^2] \}$} \\
\hline
\end{tabularx}

\newpage
% =============================================================================
% ========================= TOÁN 6 - TUẦN 05 - BUỔI 02: MÃ ĐỀ 101 =============
% ========================= MẶT SAU (TRANG 2/2) ================================
% =============================================================================

\noindent
\begin{tabularx}{\linewidth}{|X|c|}
\hline
\textbf{PHẦN II. TỰ LUẬN (Tiếp theo) - TOÁN 6 TUẦN 05 BUỔI 02} & \textbf{Mã đề: 101 - Trang 2/2} \\
\hline
\end{tabularx}

\vspace{2.5mm}

\noindent\textbf{\large Câu 2 (2,5 điểm). Tìm số tự nhiên $x$, biết:}
\vspace{0.8mm}

\noindent
\begin{tabularx}{\linewidth}{|X|X|}
\hline
\oBaiTap[6]{a)}{$24 + 3(x - 2) = 45$} & 
\oBaiTap[6]{b)}{$2 \cdot 3^x - 14 = 40$} \\
\hline
\oBaiTap[6]{c)}{Tìm số tự nhiên $x$ lớn nhất, biết rằng $48 \ \vdots \ x$ và $72 \ \vdots \ x$.} & 
\oBaiTap[6]{d)}{Tìm số tự nhiên $x$, biết rằng $60 \ \vdots \ x$, $84 \ \vdots \ x$ và $x > 6$.} \\
\hline
\end{tabularx}

\vspace{2.5mm}

\noindent\textbf{\large Câu 3 (2,5 điểm). Bài toán thực tế:}
\vspace{0.8mm}

\noindent
\begin{tabularx}{\linewidth}{|X|}
\hline
\vspace{0.5mm}
Bác Hùng có một khu đất hình chữ nhật có chiều dài là $18\text{ m}$ và chiều rộng là $12\text{ m}$. Bác dự định làm hàng rào lưới thép B40 xung quanh khu đất và để lại một cổng ra vào rộng $2\text{ m}$.\par
\textbf{a) (1,25 điểm)} Tính chu vi của khu đất hình chữ nhật đó.\par
\textbf{b) (1,25 điểm)} Tính chiều dài hàng rào lưới thép B40 cần làm và số tiền bác Hùng phải trả để mua lưới thép, biết giá mỗi mét lưới thép là $80\,000\text{ đồng}$.\par
\vspace{0.4mm}
\noindent\textit{Bài giải:}\dongke{14}
\\ \hline
\end{tabularx}

\newpage

% =============================================================================
% ========================= TOÁN 6 - TUẦN 05 - BUỔI 02: MÃ ĐỀ 102 =============
% ========================= MẶT TRƯỚC (TRANG 1/2) ==============================
% =============================================================================

\noindent
\begin{tabularx}{\linewidth}{|X|c|}
\hline
\textbf{PHIẾU KIỂM TRA ĐỊNH KỲ TOÁN 6 - TUẦN 05 - BUỔI 02} & \textbf{Mã đề: 102 - Trang 1/2} \\
\textit{Chủ đề: Dấu hiệu chia hết (2, 3, 5, 9) --- Tìm ƯCLN --- Tứ giác thực tế} & \textit{Thời gian: 45 phút} \\
\hline
\end{tabularx}

\vspace{1mm}

\noindent
\begin{tabularx}{\linewidth}{|X|}
\hline
\rule[-3.5mm]{0pt}{10mm}\textbf{Họ và tên học sinh:} \hfill \textbf{Mã HS:} \makebox[3.5cm]{} \\
\hline
\end{tabularx}

\vspace{1mm}

\noindent
\begin{tabularx}{\linewidth}{|p{3.2cm}|X|}
\hline
\centering\textbf{ĐIỂM SỐ} & \centering\textbf{LỜI NHẬN XÉT CỦA THẦY CÔ} \tabularnewline
\hline
\rule{0pt}{1.3cm} & \tabularnewline
\hline
\end{tabularx}

\vspace{1.2mm}

\noindent\textbf{BẢNG ĐIỀN ĐÁP ÁN TRẮC NGHIỆM (Mỗi câu đúng 0,5 điểm):}
\vspace{0.8mm}

\noindent
\begingroup
\setlength{\tabcolsep}{3pt}
\renewcommand{\arraystretch}{1.35}
\begin{tabularx}{\linewidth}{|c|*{5}{>{\centering\arraybackslash}X|}}
\hline
\textbf{Câu} & \textbf{1} & \textbf{2} & \textbf{3} & \textbf{4} & \textbf{5} \\
\hline
\textbf{Đáp án} & & & & & \\
\hline
\end{tabularx}
\endgroup

\vspace{1.5mm}

\noindent\textbf{\large PHẦN I. TRẮC NGHIỆM KHÁCH QUAN (2,5 điểm)}
\vspace{1mm}

\cauTN{1}{Số nào sau đây chia hết cho $2$ nhưng \textbf{không chia hết} cho $5$?}{$150$}{$274$}{$385$}{$411$}

\cauTN{2}{Số nào sau đây chia hết cho cả $3$ và $9$?}{$123$}{$234$}{$345$}{$456$}

\cauTN{3}{Cho số $\overline{4*7}$. Chữ số $*$ thích hợp để số đó chia hết cho $3$ có thể là:}{$1$}{$2$}{$3$}{$5$}

\cauTN{4}{Tổng nào sau đây chia hết cho $7$?}{$14 + 21 + 35$}{$14 + 21 + 36$}{$7 + 15 + 28$}{$21 + 30 + 49$}

\cauTNB{5}{Một số tự nhiên có chữ số tận cùng bằng $0$ và tổng các chữ số chia hết cho $3$ thì số đó:}{Chỉ chia hết cho 2 và 5}{Chỉ chia hết cho 3}{Chia hết cho cả 2, 3 và 5}{Chia hết cho cả 2, 3, 5 và 9}

\vspace{1.2mm}

\noindent\textbf{\large PHẦN II. TỰ LUẬN (7,5 điểm)}

\noindent\textbf{Câu 1 (2,5 điểm). Thực hiện phép tính (hợp lý nếu có thể):}
\vspace{0.8mm}

\noindent
\begin{tabularx}{\linewidth}{|X|X|}
\hline
\oBaiTap[6]{a)}{$42 + 25 \cdot 2 - 3^3$} & 
\oBaiTap[6]{b)}{$37 \cdot 64 + 37 \cdot 36 - 700$} \\
\hline
\oBaiTap[6]{c)}{$140 : [54 - (50 - 2^4)]$} & 
\oBaiTap[6]{d)}{$320 : \{ 180 - [75 + (11 - 6)^2] \}$} \\
\hline
\end{tabularx}

\newpage
% =============================================================================
% ========================= TOÁN 6 - TUẦN 05 - BUỔI 02: MÃ ĐỀ 102 =============
% ========================= MẶT SAU (TRANG 2/2) ================================
% =============================================================================

\noindent
\begin{tabularx}{\linewidth}{|X|c|}
\hline
\textbf{PHẦN II. TỰ LUẬN (Tiếp theo) - TOÁN 6 TUẦN 05 BUỔI 02} & \textbf{Mã đề: 102 - Trang 2/2} \\
\hline
\end{tabularx}

\vspace{2.5mm}

\noindent\textbf{\large Câu 2 (2,5 điểm). Tìm số tự nhiên $x$, biết:}
\vspace{0.8mm}

\noindent
\begin{tabularx}{\linewidth}{|X|X|}
\hline
\oBaiTap[6]{a)}{$35 + 4(x - 3) = 67$} & 
\oBaiTap[6]{b)}{$3 \cdot 2^x - 7 = 41$} \\
\hline
\oBaiTap[6]{c)}{Tìm số tự nhiên $x$ lớn nhất, biết rằng $56 \ \vdots \ x$ và $84 \ \vdots \ x$.} & 
\oBaiTap[6]{d)}{Tìm số tự nhiên $x$, biết rằng $70 \ \vdots \ x$, $105 \ \vdots \ x$ và $x > 10$.} \\
\hline
\end{tabularx}

\vspace{2.5mm}

\noindent\textbf{\large Câu 3 (2,5 điểm). Bài toán thực tế:}
\vspace{0.8mm}

\noindent
\begin{tabularx}{\linewidth}{|X|}
\hline
\vspace{0.5mm}
Bác Minh có một mảnh vườn trồng hoa hình chữ nhật có chiều dài là $20\text{ m}$ và chiều rộng là $15\text{ m}$. Bác dự định làm hàng rào dây thép gai xung quanh vườn và chừa lại một lối đi rộng $3\text{ m}$.\par
\textbf{a) (1,25 điểm)} Tính chu vi của mảnh vườn hình chữ nhật đó.\par
\textbf{b) (1,25 điểm)} Tính chiều dài dây thép gai cần làm và số tiền bác Minh phải trả, biết giá mỗi mét dây thép gai là $60\,000\text{ đồng}$.\par
\vspace{0.4mm}
\noindent\textit{Bài giải:}\dongke{14}
\\ \hline
\end{tabularx}

\newpage

% =============================================================================
% ========================= TOÁN 6 - TUẦN 05 - BUỔI 02: MÃ ĐỀ 103 =============
% ========================= MẶT TRƯỚC (TRANG 1/2) ==============================
% =============================================================================

\noindent
\begin{tabularx}{\linewidth}{|X|c|}
\hline
\textbf{PHIẾU KIỂM TRA ĐỊNH KỲ TOÁN 6 - TUẦN 05 - BUỔI 02} & \textbf{Mã đề: 103 - Trang 1/2} \\
\textit{Chủ đề: Dấu hiệu chia hết (2, 3, 5, 9) --- Tìm ƯCLN --- Tứ giác thực tế} & \textit{Thời gian: 45 phút} \\
\hline
\end{tabularx}

\vspace{1mm}

\noindent
\begin{tabularx}{\linewidth}{|X|}
\hline
\rule[-3.5mm]{0pt}{10mm}\textbf{Họ và tên học sinh:} \hfill \textbf{Mã HS:} \makebox[3.5cm]{} \\
\hline
\end{tabularx}

\vspace{1mm}

\noindent
\begin{tabularx}{\linewidth}{|p{3.2cm}|X|}
\hline
\centering\textbf{ĐIỂM SỐ} & \centering\textbf{LỜI NHẬN XÉT CỦA THẦY CÔ} \tabularnewline
\hline
\rule{0pt}{1.3cm} & \tabularnewline
\hline
\end{tabularx}

\vspace{1.2mm}

\noindent\textbf{BẢNG ĐIỀN ĐÁP ÁN TRẮC NGHIỆM (Mỗi câu đúng 0,5 điểm):}
\vspace{0.8mm}

\noindent
\begingroup
\setlength{\tabcolsep}{3pt}
\renewcommand{\arraystretch}{1.35}
\begin{tabularx}{\linewidth}{|c|*{5}{>{\centering\arraybackslash}X|}}
\hline
\textbf{Câu} & \textbf{1} & \textbf{2} & \textbf{3} & \textbf{4} & \textbf{5} \\
\hline
\textbf{Đáp án} & & & & & \\
\hline
\end{tabularx}
\endgroup

\vspace{1.5mm}

\noindent\textbf{\large PHẦN I. TRẮC NGHIỆM KHÁCH QUAN (2,5 điểm)}
\vspace{1mm}

\cauTN{1}{Số nào sau đây chia hết cho $5$ nhưng \textbf{không chia hết} cho $2$?}{$240$}{$375$}{$468$}{$512$}

\cauTN{2}{Số nào sau đây \textbf{không chia hết} cho $3$?}{$132$}{$246$}{$357$}{$418$}

\cauTNB{3}{Chữ số thích hợp điền vào dấu $*$ để số $\overline{58*}$ chia hết cho cả $2$ và $5$ là:}{$0$}{$2$}{$5$}{Không có chữ số nào}

\cauTNB{4}{Cho tổng $S = 24 + 36 + x$ (với $x \in \mathbb{N}$). Để $S$ chia hết cho $6$ thì:}{$x$ là số tự nhiên bất kỳ}{$x$ phải chia hết cho 6}{$x$ không chia hết cho 6}{$x$ chia cho 6 dư 1}

\cauTN{5}{Trong các số sau: $342, 516, 725, 814$, số chia hết cho $9$ là:}{$342$}{$516$}{$725$}{$814$}

\vspace{1.2mm}

\noindent\textbf{\large PHẦN II. TỰ LUẬN (7,5 điểm)}

\noindent\textbf{Câu 1 (2,5 điểm). Thực hiện phép tính (hợp lý nếu có thể):}
\vspace{0.8mm}

\noindent
\begin{tabularx}{\linewidth}{|X|X|}
\hline
\oBaiTap[6]{a)}{$45 + 18 \cdot 3 - 2^5$} & 
\oBaiTap[6]{b)}{$52 \cdot 73 + 52 \cdot 27 - 1200$} \\
\hline
\oBaiTap[6]{c)}{$180 : [60 - (46 - 2^4)]$} & 
\oBaiTap[6]{d)}{$400 : \{ 210 - [85 + (9 - 4)^2] \}$} \\
\hline
\end{tabularx}

\newpage
% =============================================================================
% ========================= TOÁN 6 - TUẦN 05 - BUỔI 02: MÃ ĐỀ 103 =============
% ========================= MẶT SAU (TRANG 2/2) ================================
% =============================================================================

\noindent
\begin{tabularx}{\linewidth}{|X|c|}
\hline
\textbf{PHẦN II. TỰ LUẬN (Tiếp theo) - TOÁN 6 TUẦN 05 BUỔI 02} & \textbf{Mã đề: 103 - Trang 2/2} \\
\hline
\end{tabularx}

\vspace{2.5mm}

\noindent\textbf{\large Câu 2 (2,5 điểm). Tìm số tự nhiên $x$, biết:}
\vspace{0.8mm}

\noindent
\begin{tabularx}{\linewidth}{|X|X|}
\hline
\oBaiTap[6]{a)}{$28 + 5(x - 4) = 73$} & 
\oBaiTap[6]{b)}{$5 \cdot 2^x - 12 = 68$} \\
\hline
\oBaiTap[6]{c)}{Tìm số tự nhiên $x$ lớn nhất, biết rằng $54 \ \vdots \ x$ và $90 \ \vdots \ x$.} & 
\oBaiTap[6]{d)}{Tìm số tự nhiên $x$, biết rằng $45 \ \vdots \ x$, $75 \ \vdots \ x$ và $x > 5$.} \\
\hline
\end{tabularx}

\vspace{2.5mm}

\noindent\textbf{\large Câu 3 (2,5 điểm). Bài toán thực tế:}
\vspace{0.8mm}

\noindent
\begin{tabularx}{\linewidth}{|X|}
\hline
\vspace{0.5mm}
Một khung cửa sổ bằng nhôm kính hình chữ nhật có chiều dài là $160\text{ cm}$ và chiều rộng là $120\text{ cm}$. Người thợ cần lắp đặt dải gioăng cao su bảo vệ bao quanh toàn bộ viền ngoài của khung cửa sổ.\par
\textbf{a) (1,25 điểm)} Tính chu vi của khung cửa sổ đó (theo đơn vị mét).\par
\textbf{b) (1,25 điểm)} Tính tổng chi phí mua dải gioăng cao su, biết đơn giá mỗi mét gioăng là $45\,000\text{ đồng}$.\par
\vspace{0.4mm}
\noindent\textit{Bài giải:}\dongke{14}
\\ \hline
\end{tabularx}

\newpage

% =============================================================================
% ========================= TOÁN 6 - TUẦN 05 - BUỔI 02: MÃ ĐỀ 104 =============
% ========================= MẶT TRƯỚC (TRANG 1/2) ==============================
% =============================================================================

\noindent
\begin{tabularx}{\linewidth}{|X|c|}
\hline
\textbf{PHIẾU KIỂM TRA ĐỊNH KỲ TOÁN 6 - TUẦN 05 - BUỔI 02} & \textbf{Mã đề: 104 - Trang 1/2} \\
\textit{Chủ đề: Dấu hiệu chia hết (2, 3, 5, 9) --- Tìm ƯCLN --- Tứ giác thực tế} & \textit{Thời gian: 45 phút} \\
\hline
\end{tabularx}

\vspace{1mm}

\noindent
\begin{tabularx}{\linewidth}{|X|}
\hline
\rule[-3.5mm]{0pt}{10mm}\textbf{Họ và tên học sinh:} \hfill \textbf{Mã HS:} \makebox[3.5cm]{} \\
\hline
\end{tabularx}

\vspace{1mm}

\noindent
\begin{tabularx}{\linewidth}{|p{3.2cm}|X|}
\hline
\centering\textbf{ĐIỂM SỐ} & \centering\textbf{LỜI NHẬN XÉT CỦA THẦY CÔ} \tabularnewline
\hline
\rule{0pt}{1.3cm} & \tabularnewline
\hline
\end{tabularx}

\vspace{1.2mm}

\noindent\textbf{BẢNG ĐIỀN ĐÁP ÁN TRẮC NGHIỆM (Mỗi câu đúng 0,5 điểm):}
\vspace{0.8mm}

\noindent
\begingroup
\setlength{\tabcolsep}{3pt}
\renewcommand{\arraystretch}{1.35}
\begin{tabularx}{\linewidth}{|c|*{5}{>{\centering\arraybackslash}X|}}
\hline
\textbf{Câu} & \textbf{1} & \textbf{2} & \textbf{3} & \textbf{4} & \textbf{5} \\
\hline
\textbf{Đáp án} & & & & & \\
\hline
\end{tabularx}
\endgroup

\vspace{1.5mm}

\noindent\textbf{\large PHẦN I. TRẮC NGHIỆM KHÁCH QUAN (2,5 điểm)}
\vspace{1mm}

\cauTN{1}{Các số có chữ số tận cùng là $0, 2, 4, 6, 8$ thì chắc chắn:}{Chia hết cho 5}{Chia hết cho 3}{Chia hết cho 2}{Chia hết cho cả 2 và 5}

\cauTNB{2}{Khẳng định nào sau đây là \textbf{đúng}?}{Mọi số chia hết cho 3 đều chia hết cho 9}{Mọi số chia hết cho 9 đều chia hết cho 3}{Số có chữ số tận cùng là 9 thì chia hết cho 9}{Số có tổng các chữ số bằng 6 thì chia hết cho 9}

\cauTN{3}{Chữ số $*$ thích hợp để số $\overline{7*5}$ chia hết cho $9$ là:}{$3$}{$5$}{$6$}{$9$}

\cauTN{4}{Hiệu nào sau đây chia hết cho $8$?}{$64 - 24$}{$56 - 25$}{$48 - 18$}{$72 - 30$}

\cauTN{5}{Số tự nhiên nhỏ nhất có ba chữ số khác nhau chia hết cho cả $2$ và $5$ là:}{$100$}{$120$}{$102$}{$105$}

\vspace{1.2mm}

\noindent\textbf{\large PHẦN II. TỰ LUẬN (7,5 điểm)}

\noindent\textbf{Câu 1 (2,5 điểm). Thực hiện phép tính (hợp lý nếu có thể):}
\vspace{0.8mm}

\noindent
\begin{tabularx}{\linewidth}{|X|X|}
\hline
\oBaiTap[6]{a)}{$56 + 15 \cdot 3 - 3^3$} & 
\oBaiTap[6]{b)}{$64 \cdot 82 + 64 \cdot 18 - 1400$} \\
\hline
\oBaiTap[6]{c)}{$200 : [55 - (46 - 2^4)]$} & 
\oBaiTap[6]{d)}{$500 : \{ 220 - [95 + (8 - 3)^2] \}$} \\
\hline
\end{tabularx}

\newpage
% =============================================================================
% ========================= TOÁN 6 - TUẦN 05 - BUỔI 02: MÃ ĐỀ 104 =============
% ========================= MẶT SAU (TRANG 2/2) ================================
% =============================================================================

\noindent
\begin{tabularx}{\linewidth}{|X|c|}
\hline
\textbf{PHẦN II. TỰ LUẬN (Tiếp theo) - TOÁN 6 TUẦN 05 BUỔI 02} & \textbf{Mã đề: 104 - Trang 2/2} \\
\hline
\end{tabularx}

\vspace{2.5mm}

\noindent\textbf{\large Câu 2 (2,5 điểm). Tìm số tự nhiên $x$, biết:}
\vspace{0.8mm}

\noindent
\begin{tabularx}{\linewidth}{|X|X|}
\hline
\oBaiTap[6]{a)}{$36 + 6(x - 5) = 84$} & 
\oBaiTap[6]{b)}{$4 \cdot 3^x - 22 = 86$} \\
\hline
\oBaiTap[6]{c)}{Tìm số tự nhiên $x$ lớn nhất, biết rằng $64 \ \vdots \ x$ và $96 \ \vdots \ x$.} & 
\oBaiTap[6]{d)}{Tìm số tự nhiên $x$, biết rằng $40 \ \vdots \ x$, $60 \ \vdots \ x$ và $x > 10$.} \\
\hline
\end{tabularx}

\vspace{2.5mm}

\noindent\textbf{\large Câu 3 (2,5 điểm). Bài toán thực tế:}
\vspace{0.8mm}

\noindent
\begin{tabularx}{\linewidth}{|X|}
\hline
\vspace{0.5mm}
Một biển hiệu quảng cáo có dạng hình thoi với độ dài mỗi cạnh là $2{,}5\text{ m}$. Ban quản lý cửa hàng cần lắp dải đèn LED trang trí chạy dọc theo toàn bộ viền ngoài của biển quảng cáo đó.\par
\textbf{a) (1,25 điểm)} Tính chu vi của biển hiệu quảng cáo hình thoi đó.\par
\textbf{b) (1,25 điểm)} Tính tổng số tiền cần thanh toán để mua dải đèn LED, biết giá tiền mỗi mét đèn LED là $115\,000\text{ đồng}$.\par
\vspace{0.4mm}
\noindent\textit{Bài giải:}\dongke{14}
\\ \hline
\end{tabularx}

\end{document}
